\origpage{69}{GRAPH MINORS. XIII}
For example, let $G$ be the digraph with four vertices $u, v_1, v_2, v_3$ and three edges $uv_1, uv_2, uv_3$; we shall examine the $0$-folio of $(G, v_1, v_2, v_3)$. For $i = 1, 2, 3$ let $N_i$ be the digraph with vertex set $\{1, \dots, i\}$ and no edges. Then the $0$-folio of $(G, v_1, v_2, v_3)$ contains $(N_1, 1, 1, 1)$, $(N_2, 1, 1, 2)$, $(N_2, 1, 2, 1)$, $(N_2, 2, 1, 1)$, $(N_3, 1, 2, 3)$, and all rooted digraphs isomorphic to one of these.

Now the $\delta$-folio of $(G, v_1, \dots, v_k)$ can be determined from the $\delta$-folio of $(G, u_1, \dots, u_l)$ if $\{u_1, \dots, u_l\} = \{v_1, \dots, v_k\}$, as is easily seen. Thus, if $Z \subseteq V(G)$, by the $\delta$-folio of $G$ relative to $Z$ we mean the $\delta$-folio of $(G, v_1, \dots, v_k)$ for any choice of $v_1, \dots, v_k$ with $k = |Z|$ and $Z = \{v_1, \dots, v_k\}$.

For fixed $\delta, \xi \ge 0$ we shall give an algorithm with running time $O(v^3)$ to solve the following problem.

\textit{Folio.} Instance: A digraph $G$, and a subset $Z \subseteq V(G)$ with $|Z| \le \xi$. Question: What is the $\delta$-folio of $G$ relative to $Z$?

Let us see how Folio can be used to solve the two problems of the previous section. For Disjoint Connected Subgraphs, with input $G = (V, E)$ and $Z_1, \dots, Z_t \subseteq V$, we may assume that $Z_1, \dots, Z_t$ are mutually disjoint; we direct the edges of $G$ arbitrarily; we compute the $0$-folio of $G$ relative to $Z_1 \cup \dots \cup Z_t$, and from this we read off whether the desired connected subgraphs exist. (Compare the example above.) For $H$-Minor Containment, we direct the edges of $G$ arbitrarily and compute the $\delta$-folio of $G$ relative to $\emptyset$, where $\delta = \max(|V(H)|, |E(H)|)$. These both have running time $O(v^3)$.

We shall need some easy observations on folio management\del{, the proofs of which we leave to the reader}:

\stmt
(2.1) \textit{If $Z' \subseteq Z \subseteq V(G)$ and $\delta \ge 0$, the $\delta$-folio of $G$ relative to $Z'$ is determined by a knowledge of the $\delta$-folio of $G$ relative to $Z$.}

\fix{\textit{Proof of (2.1).} Let $k = |Z|$ and $k' = |Z'|$, and let $Z = \{v_1, \dots, v_k\}$, where $Z' = \{v_1, \dots, v_{k'}\}$. Let $(H, u_1, \dots, u_{k'})$ be a rooted digraph with detail $\le \delta$. We say that a rooted digraph $(H^+, u_1, \dots, u_k)$ \textit{extends} it if $H$ is a subdigraph of $H^+$ with $E(H^+) = E(H)$, and every vertex of $H^+$ not in $V(H)$ is one of $u_{k'+1}, \dots, u_k$. Such a rooted digraph has detail $\le \delta$, because it has the same edges as $H$, and each of its vertices that is not a root is a vertex of $H$ different from $u_1, \dots, u_{k'}$. We claim that $(H, u_1, \dots, u_{k'})$ is a minor of $(G, v_1, \dots, v_{k'})$ if and only if some rooted digraph that extends it is a minor of $(G, v_1, \dots, v_k)$. Then the $\delta$-folio of $(G, v_1, \dots, v_{k'})$ consists of the rooted digraphs with detail $\le \delta$ that are extended by some member of the $\delta$-folio of $(G, v_1, \dots, v_k)$, and (2.1) follows.

To prove the claim, let $\phi$ be a model in $(G, v_1, \dots, v_k)$ of a rooted digraph that extends $(H, u_1, \dots, u_{k'})$. Then the restriction of $\phi$ to $V(H) \cup E(H)$ is a model of $(H, u_1, \dots, u_{k'})$ in $(G, v_1, \dots, v_{k'})$. Conversely, let $\phi$ be a model of $(H, u_1, \dots, u_{k'})$ in $(G, v_1, \dots, v_{k'})$. For $k' < i \le k$, if $v_i \in V(\phi(u))$ for some $u \in V(H)$ let $u_i = u$, and otherwise let $u_i$ be a new vertex and let $\phi(u_i)$ be the subgraph of $G$ with vertex set $\{v_i\}$ and no edges. Let $H^+$ be obtained from $H$ by adding the new vertices. Then $(H^+, u_1, \dots, u_k)$ extends $(H, u_1, \dots, u_{k'})$, and $\phi$ is a model of it in $(G, v_1, \dots, v_k)$, because the vertices $v_i$ with $u_i$ new are distinct and belong to no $\phi(u)$ with $u \in V(H)$. This proves (2.1). $\square$}\fixnote{fix:2.1-proof-added}{The paper leaves the proof of (2.1) to the reader. The argument above is the one used in the formalization.}

A separation of a graph or digraph $G$ is a pair $(A, B)$ of subgraphs such that $A \cup B = G$ and $E(A \cap B) = \emptyset$, where for subgraphs $G_1 = (V_1, E_1)$, $G_2 = (V_2, E_2)$ of $G$, $G_1 \cup G_2$ and $G_1 \cap G_2$ are the subgraphs $(V_1 \cup V_2, E_1 \cup E_2)$ and $(V_1 \cap V_2, E_1 \cap E_2)$, respectively. The order of the separation $(A, B)$ is $|V(A \cap B)|$.

\stmt
(2.2) \textit{Let $(A, B)$ be a separation of a digraph $G$, and let $V(A \cap B) \subseteq Z \subseteq V(G)$. For $\delta \ge 0$, the $\delta$-folio of $G$ relative to $Z$ is determined by a knowledge of the $\delta$-folios of $A$ relative to $Z \cap V(A)$ and of $B$ relative to $Z \cap V(B)$.}

\fix{\textit{Proof of (2.2).} Let $Z = \{v_1, \dots, v_k\}$ with $k = |Z|$. Let $I_A$ be the set of all $i$ ($1 \le i \le k$) with $v_i \in V(A)$, and define $I_B$ similarly. Thus every $i$ belongs to $I_A$ or to $I_B$, and $i \in I_A \cap I_B$ if and only if $v_i \in V(A \cap B)$. Let $(H, u_1, \dots, u_k)$ be a rooted digraph with detail $\le \delta$. A \textit{split} of it consists of rooted digraphs $(H_A, a_i : i \in I_A)$ and $(H_B, b_i : i \in I_B)$ with roots indexed by $I_A$ and $I_B$, and maps $\pi_A \colon V(H_A) \to V(H)$ and $\pi_B \colon V(H_B) \to V(H)$, such that

\hspace*{12pt}(a) $E(H)$ is the disjoint union of $E(H_A)$ and $E(H_B)$, and each $e \in E(H_A)$ with head $a$ and tail $a'$ in $H_A$ has head $\pi_A(a)$ and tail $\pi_A(a')$ in $H$, and similarly for $H_B$

\hspace*{12pt}(b) $\pi_A(a_i) = u_i$ for $i \in I_A$, and $\pi_B(b_i) = u_i$ for $i \in I_B$

\hspace*{12pt}(c) for each $u \in V(H)$ the set $\pi_A^{-1}(u) \cup \pi_B^{-1}(u)$ is non-empty, and the graph $K_u$ with this vertex set, in which $a_i$ and $b_i$ are adjacent for each $i \in I_A \cap I_B$ with $u_i = u$, is connected.

We claim that $(H, u_1, \dots, u_k)$ is a minor of $(G, v_1, \dots, v_k)$ if and only if it has a split in which $(H_A, a_i : i \in I_A)$ and $(H_B, b_i : i \in I_B)$ have detail $\le \delta$ and are minors of $(A, v_i : i \in I_A)$ and $(B, v_i : i \in I_B)$ respectively. This condition depends only on the $\delta$-folios of $A$ relative to $Z \cap V(A)$ and of $B$ relative to $Z \cap V(B)$, and so (2.2) follows from the claim.

Suppose first that there is such a split, and let $\phi_A, \phi_B$ be models of the two rooted digraphs in $(A, v_i : i \in I_A)$ and $(B, v_i : i \in I_B)$. For $u \in V(H)$ let $\phi(u)$ be the union of the graphs $\phi_A(a)$ ($a \in \pi_A^{-1}(u)$) and $\phi_B(b)$ ($b \in \pi_B^{-1}(u)$), and for $e \in E(H)$ let $\phi(e) = \phi_A(e)$ if $e \in E(H_A)$ and $\phi(e) = \phi_B(e)$ if $e \in E(H_B)$. By (c), $\phi(u)$ is non-null. It is connected, since its parts are connected, $K_u$ is connected, and $v_i \in V(\phi_A(a_i)) \cap V(\phi_B(b_i))$ for $i \in I_A \cap I_B$. Let $u, w \in V(H)$ be distinct, and suppose that some $x$ belongs to $V(\phi(u)) \cap V(\phi(w))$. Since the graphs $\phi_A(a)$ ($a \in V(H_A)$) are mutually disjoint, and so are the graphs $\phi_B(b)$ ($b \in V(H_B)$), we may assume that $x \in V(\phi_A(a)) \cap V(\phi_B(b))$ where $\pi_A(a) = u$ and $\pi_B(b) = w$. Then $x \in V(A \cap B) \subseteq Z$, and so $x = v_i$ for some $i \in I_A \cap I_B$. Since $v_i \in V(\phi_A(a_i))$, it follows that $a = a_i$, and similarly $b = b_i$. So $u = u_i = w$ by (b), a contradiction. Thus the graphs $\phi(u)$ are mutually vertex-disjoint. The edges $\phi(e)$ are distinct, and none of them is an edge of any $\phi(u)$, because $\phi_A$ and $\phi_B$ are models and $E(A \cap B) = \emptyset$. By (a), each $\phi(e)$ has its head and tail in the right graphs $\phi(u)$, and $v_i \in V(\phi(u_i))$ for $1 \le i \le k$ by (b). Hence $\phi$ is a model of $(H, u_1, \dots, u_k)$ in $(G, v_1, \dots, v_k)$.

Conversely, let $\phi$ be a model of $(H, u_1, \dots, u_k)$ in $(G, v_1, \dots, v_k)$. For $u \in V(H)$, we call the components of the subgraph $(V(\phi(u)) \cap V(A), E(\phi(u)) \cap E(A))$ the \textit{$A$-pieces} of $u$, and define the $B$-pieces of $u$ similarly. Let $H_A$ be the digraph whose vertices are the $A$-pieces of all vertices of $H$, and whose edges are the edges $e \in E(H)$ with $\phi(e) \in E(A)$, where such an $e$ has as head the $A$-piece containing the head of $\phi(e)$ and as tail the $A$-piece containing the tail of $\phi(e)$. (If $u$ is the head of $e$ in $H$, the head of $\phi(e)$ is in $V(\phi(u)) \cap V(A)$, and so it is in an $A$-piece of $u$.) For $i \in I_A$ let $a_i$ be the $A$-piece of $u_i$ that contains $v_i$, and let $\pi_A$ map each $A$-piece of $u$ to $u$. Define $H_B$, $b_i$ and $\pi_B$ similarly. Mapping each piece to itself and each edge $e$ to $\phi(e)$ gives models of $(H_A, a_i : i \in I_A)$ in $(A, v_i : i \in I_A)$ and of $(H_B, b_i : i \in I_B)$ in $(B, v_i : i \in I_B)$, and (a) and (b) hold by construction. For (c), let $u \in V(H)$. Then $\phi(u)$ is non-null, and so $u$ has a piece. If an $A$-piece and a $B$-piece of $u$ have a common vertex $x$, then $x \in V(A \cap B) \subseteq Z$, and so $x = v_i$ for some $i \in I_A \cap I_B$. Then $u = u_i$, since $v_i \in V(\phi(u_i))$, and the two pieces are $a_i$ and $b_i$, which are adjacent in $K_u$. Since every edge of $\phi(u)$ is an edge of one piece, no two $A$-pieces and no two $B$-pieces of $u$ have a common vertex, and $\phi(u)$ is connected, it follows that $K_u$ is connected. It remains to show that $(H_A, a_i : i \in I_A)$ has detail $\le \delta$ (and similarly for $H_B$). First, $|E(H_A)| \le |E(H)| \le \delta$. Now let $P$ be an $A$-piece of some $u \in V(H)$ that is not a root of $H_A$. Then $V(P) \cap Z = \emptyset$, for if $v_i \in V(P)$ then $i \in I_A$, $u = u_i$ and $P = a_i$. Let $f$ be an edge of $\phi(u)$ with an end $x \in V(P)$. If $f \in E(B)$ then $x \in V(A \cap B) \subseteq Z$, which is impossible. So $f \in E(A)$, and $f \in E(P)$. Since $\phi(u)$ is connected, it follows that $P = \phi(u)$. Moreover $u$ is not a root of $H$, for if $u = u_i$ then $v_i \in V(\phi(u)) = V(P)$. Hence the vertices of $H_A$ that are not roots are graphs $\phi(u)$ for distinct vertices $u$ of $H$ that are not roots, and there are at most $\delta$ of them. This proves the claim, and hence (2.2). $\square$}\fixnote{fix:2.2-proof-added}{The paper leaves the proof of (2.2) to the reader. The argument above, which cuts a model along the separation into pieces and glues models of the pieces, is the one used in the formalization.}

For $Z \subseteq V(G)$, a vertex $v$ of $G$ is \textit{irrelevant (to the $\delta$-folio of $G$ relative to $Z$)} if $v \notin Z$ and the $\delta$-folios of $G$ and $G \setminus v$ relative to $Z$ are equal. (We shall use $\setminus$ for deletion.) From (2.1) and (2.2), the following follows easily.
